\paragraph{ConvLSTM.}
This baseline adapts the CATCH architecture~\citep{agmata2025catch}, proposed for forecasting fisheries catch hotspots, to the prediction of Poisson counts. It is built on convolutional LSTM (ConvLSTM) cells~\citep{shi2015convolutional}, which replace the matrix products of an LSTM with convolutions, so that hidden and cell states are spatial maps. The grid cells are arranged on the smallest regular raster that contains them ($7\times14$ positions), and a fixed binary mask restricts predictions and the loss to the cells of the study area. To predict day $t$, the network processes the frames of the previous $T=3$ days, each with one channel per covariate and one channel with the observed counts, through two stacked ConvLSTM layers with four hidden channels, $3\times3$ convolutions and dropout of $0.2$ after each layer. The final hidden state is concatenated with the covariate frame of day $t$, and a $3\times3$ convolutional output head followed by a softplus activation returns a positive Poisson rate $\widehat{\lambda}_{j,t}$ for every cell. All input channels are standardized with training statistics, and the network is trained by minimizing the mean Poisson negative log-likelihood over the valid cells of all training days, using full-batch Adam with learning rate $10^{-2}$ for 300 epochs; $T$ and the learning rate were selected on the validation weeks (App.~\ref{app:hyperparameters}). During evaluation, the input frames of each test day contain the counts observed on the preceding days.
